% Compile with: latexmk -cd -pdf -interaction=nonstopmode -outdir=. spanish_v.tex
\documentclass[conference]{IEEEtran}

\usepackage{amsmath, amssymb} % Útil para ecuaciones
\usepackage{graphicx} % Para incluir imágenes
\usepackage{cite} % Para gestión de bibliografías
\usepackage{array} % Para el control de ancho de columna p{}
\usepackage{xcolor}
\usepackage{tcolorbox}

\date{\today}

\begin{document}


\title{Hacia un Marco Práctico de Protocolos Criptográficos en Cooperación Adversarial}
\author{
    \IEEEauthorblockN{S. Restrepo Ruiz}
    \IEEEauthorblockA{waajacu.com \\ Correo: santiago.restrepo.ruiz@gmail.com}
}

\maketitle

\begin{abstract}
Este estudio investiga el papel de la teoría de la información en el diseño de protocolos criptográficos que establecen confianza en entornos adversariales. 
Aprovechando computaciones seguras multi-parte, así como privacidad diferencial y pruebas de conocimiento cero, la investigación busca construir 
una biblioteca de software \textit{de código abierto} que demuestre resultados de cooperación.
\end{abstract}

\vspace{-5pt}
\section{Introducción}
Si se concede el permiso, en colaboración con la \textit{Academia China de Ciencias}, se busca desarrollar marcos criptográficos que mejoren la cooperación en interacciones estructuradas. 
Se implementarán protocolos de cifrado seguros y verificables, dirigidos a asegurar la coordinación basada en teoría de juegos y la toma de decisiones con preservación de la privacidad.

\vspace{-5pt}
\section{Revisión de la Literatura}
La criptografía va mucho más allá del cifrado, sirve como la base de la confianza digital. Las herramientas que se emplearán en este estudio son: 
\textcolor{blue}{\textit{pruebas de conocimiento cero}}, 
\textcolor{blue}{\textit{cálculo homomórfico}}, 
\textcolor{blue}{\textit{circuitos Garbled}}, 
\textcolor{blue}{\textit{transferencia oblivia}}, 
\textcolor{blue}{\textit{blockchains}}, 
\textcolor{blue}{\textit{privacidad diferencial}} y 
\textcolor{blue}{\textit{criptografía de umbral}}. 
---La combinación de estas unidades seguras posibilita la resolución colaborativa de problemas que, en mi opinión, son la clave para la cooperación adversarial.

\vspace{-5pt}
\section{Preguntas de Investigación}

\begin{tabular}{lp{210pt}}
    {\small \textbf{1)}} & {\small ¿Qué implica aplicar la criptografía a la cooperación?} \\ 
    {\small \textbf{2)}} & {\small ¿Cómo pueden mejorarse los protocolos criptográficos propuestos mediante la colaboración interdisciplinaria?}
\end{tabular}
    
\vspace{-5pt}
\section{Objetivos de la Investigación}

\begin{tabular}{lp{210pt}}
    {\small \textbf{1)}} & {\small Demostrar en problemas clásicos de teoría de juegos.} \\ 
    {\small \textbf{2)}} & {\small Publicar una biblioteca de software de código abierto basada en esta investigación.}
\end{tabular}

\vspace{-5pt}
\section{Metodología}
El estudio se centrará en los siguientes juegos estructurados:

\begin{tabular}{lp{210pt}}
    {\small \textbf{1)}} & {\textcolor{blue}{\textit{Jugar juegos de cartas sin revelar las cartas:}} \linebreak ---Demostrar cómo proclamar la victoria en un juego de cartas sin exponer las cartas en mano.} \\
    {\small \textbf{2)}} & {\textcolor{blue}{\textit{Jugar tic-tac-toe sin revelar la estrategia:}} \linebreak ---Demostrar cómo un jugador puede probar que tiene una estrategia ganadora en tic-tac-toe sin mostrar los movimientos reales.} \\
    {\small \textbf{3)}} & {\textcolor{blue}{\textit{Juego clásico de ciervos (stag-hunt) adversarial:}} \linebreak ---Demostrar un camino hacia el equilibrio para este juego fundamental de cooperación adversarial.}
\end{tabular}

\vspace{10pt}
\section{Solicitud de Financiamiento y Justificación}
Solicito humildemente el amable apoyo de los distinguidos miembros de la Embajada de China en mi proceso de solicitud. 
Tal como se indica en las directrices del Consejo Chino de Becas (www.campuschina.org) :

\begin{tcolorbox}[colback=gray!5!white, colframe=gray!10!white, sharp corners]
    \textcolor{blue}{\textit{“Los solicitantes deberán obtener la aprobación y recomendación del país de acogida antes de presentar la solicitud.”}}
\end{tcolorbox}

Las universidades en Colombia actualmente carecen de programas de investigación en este campo. He estado estudiando el idioma chino durante varios años, 
reconociendo la importancia de una comunicación y colaboración fluidas en un entorno académico chino. 
Me entusiasma profundamente la perspectiva de aprender y contribuir a la dinámica comunidad investigadora en China. 
Creo que esta beca no solo me permitiría perseguir mis aspiraciones profesionales, sino que también fomentaría un valioso intercambio académico entre nuestras naciones.

\vspace{10pt}
\section{Universidades y Programas Objetivo}
\begin{tabular}{rl}
    \textcolor{blue}{\small\textit{Universidad de Wuhan}} & {\small Maestría en Ingeniería Eléctrica}\\
    \textcolor{blue}{\small\textit{Universidad de Hubei}} & {\small Phd Matemáticas Aplicadas (Criptografía)}\\
    \textcolor{blue}{\small\textit{Universidad de Chang'an}} & {\small Maestría en Teoría de la Criptografía}\\
    \textcolor{blue}{\small\textit{Universidad de Tsinghua}} & {\small Maestría en Computación Avanzada}
\end{tabular}

\vspace{10pt}
\section{Resultados Esperados}

\begin{tabular}{lp{210pt}}
    {\small \textbf{1)}} & {\small Demostración de la cooperación adversarial en la teoría de juegos.} \\
    {\small \textbf{2)}} & {\small Desarrollo de un conjunto de protocolos criptográficos de código abierto que posibiliten la cooperación en entornos adversariales.} \\
    {\small \textbf{3)}} & {\small Abogar sobre cómo la criptografía facilita la confianza en situaciones de conflicto.}
\end{tabular}

\vspace{10pt}
\section{Conclusión}
Esta investigación busca mostrar cómo las técnicas criptográficas pueden emplearse para la construcción de paz y el establecimiento de confianza. 

% \begin{thebibliography}{99}

% \bibitem{Yao1982}
% A. C. Yao, ``Protocols for secure computations,'' in *Proceedings of the 23rd Annual Symposium on Foundations of Computer Science*, 1982.

% \bibitem{Goldwasser1989}
% S. Goldwasser, S. Micali, and C. Rackoff, ``The Knowledge Complexity of Interactive Proof Systems,'' in *SIAM Journal on Computing*, 1989.

% \bibitem{Gentry2009}
% C. Gentry, ``Fully Homomorphic Encryption Using Ideal Lattices,'' in *Proceedings of the 41st ACM Symposium on Theory of Computing (STOC)*, 2009.

% \bibitem{Dwork2006}
% C. Dwork, ``Differential Privacy,'' in *Proceedings of the International Colloquium on Automata, Languages, and Programming (ICALP)*, 2006.

% \bibitem{Canetti2000}
% R. Canetti, ``Security and Composition of Multi-Party Cryptographic Protocols,'' in *Journal of Cryptology*, 2000.

% \end{thebibliography}

% \bibliographystyle{IEEEtran}
% \bibliography{references}



\end{document}
